\begin{answers}{ch:pretraining}{事前学習の実際：データから基盤モデルまで}
\item Jaccard 係数は $3/6=0.5$ です。
「先生と書く」の 2-gram は $A=\{$先生, 生と, と書, 書く$\}$ の 4 個です。
「先生と書いた」は $B=\{$先生, 生と, と書, 書い, いた$\}$ の 5 個です。
共通するのは「先生」「生と」「と書」の 3 個で、どちらかに含まれるのは $4+5-3=6$ 個です。
式\eqref{eq:pt-jaccard}より $J(A,B)=3/6=0.5$ となります。
文末が 2 文字変わっただけで 0.5 まで下がるのは、文が短いからです。
長い文章なら、同じ違いでも係数は 1 に近いままです。
\item $t=100$ では $5.05\times10^{-4}$、$t=3050$ では約 $2.32\times10^{-4}$ です。
$t=100$ は $t_w=200$ より小さいので、式\eqref{eq:pt-lr}の上の行（ウォームアップ）を使います。
$\eta=10^{-3}\times(100+1)/200=5.05\times10^{-4}$ です。
$t=3050$ は下の行（コサイン減衰）で、進み具合は $p=(3050-200)/(4000-200)=0.75$ です。
$\cos(0.75\pi)\approx-0.707$ なので、$(1+\cos(\pi p))/2\approx0.146$ です。
$\eta=10^{-3}\times(0.1+0.9\times0.146)\approx2.32\times10^{-4}$ となります。
\item 空行を除く 27,144 行のうち、2,694 行（約 10\%）が落ちました（実測）。
内訳は、10 文字未満が 2,526 行、同じ文字が 3 割を超える行が 149 行、日本語の文字が半分以下の行が 19 行です。
落ちた行の多くは「そうです」のような短い会話や章番号で、質が悪いわけではありません。
短い会話は小説の大事な一部です。
ウェブ向けに作った規則を別の種類のデータにそのまま使うと、良い行まで捨ててしまいます。
落ちた行を必ず自分の目で見て、規則を調整してください。
URL の規則（行が \code{http://} か \code{https://} を含めば \code{False} を返す 1 行）を足しても、結果は変わりません。
漱石の作品は 1900 年代から 1910 年代に書かれたので、URL は 1 つもないからです。
\item 3,500 ステップで止めるべきだったとは言えません。
まず、ぶれです。
\file{base.pt} の検証損失を、学習中と同じ 20 バッチの平均で 10 回測り直すと、3.27 から 3.59 までばらつきました（実測。平均 3.43、標準偏差 0.11）。
3.39 から 3.48 への上昇（0.09）は、このぶれの中に収まります。
また、この検証損失は、小説（4.86）とドリル（1.31）の部分を合わせた平均です（\secref{sec:pt-base}）。
上昇がどちらの部分で起きたのかは、この平均だけでは分かりません。止めるかどうかを決めるなら、二つの部分を分けて、それぞれの推移を見る必要があります。
次に、データです。
異なるトークンは約 248 万個で、パラメータ数とほぼ同じなので、過学習が始まっていてもおかしくはありません。
ただ、最後の 500 ステップは学習率が最大値の 14\% 以下で、訓練損失も 3.12 から 3.11 とほとんど動いていません。
最後に、計算量です。
500 ステップは全体の 8 分の 1（約 4 分）で、止めても節約はわずかです。
確かめるには、毎回同じ多数のバッチで検証損失を測り、3,500 と 4,000 の両方のチェックポイントを残して比べます。
\end{answers}
